% !TEX program = lualatex
\documentclass[11pt,letterpaper]{scrartcl}
\usepackage[mode=notes,theme=graphite,fontset=lm]{synaptic}
\SynapticSetup{theme=paper}
\SynapticTitle{Field Notes on Metric Spaces}
\SynapticSubtitle{Continuity, completeness, and compactness}
\SynapticShortTitle{Metric-space notes}
\SynapticAuthor{Ada Student}
\SynapticDate{Week 3}
\SynapticTags{analysis \textbullet\ topology}
\begin{abstract}
  A working sheet for separating local, sequential, and covering arguments.
\end{abstract}

\begin{document}
\SynapticMakeTitle
\section{Metric-space toolkit}
The package respects the paper size selected by the document class.
\begin{lemma}
  Structural options are frozen after loading; themes remain switchable.
\end{lemma}
\begin{synkeyidea}
  Separate topological assumptions from metric estimates.
\end{synkeyidea}
\begin{synquestion}[For review]
  Which implications require completeness?
\end{synquestion}
\subsection{A reusable estimate}
For a contraction $T$ with constant $0<q<1$, the iterates satisfy
\[
  d(T^n x,T^n y)\leq q^n d(x,y).
\]
The estimate separates the geometric input from the limiting argument.
\begin{syntodo}[Before the exam]
  Compare sequential compactness, total boundedness, and completeness.
\end{syntodo}
\begin{synsummary}
  Definitions, results, and open questions now share a compact visual hierarchy.
\end{synsummary}
\end{document}
